\documentclass[12pt, a4paper]{article}
\usepackage[utf8]{inputenc}
\usepackage{amsmath, amssymb, amsthm, physics}
\usepackage{geometry}
\usepackage{booktabs}
\usepackage{hyperref}
​\geometry{margin=1in}
​% -------------------------------------------------------------
% BOILERPLATE / METADATA HEADER
% -------------------------------------------------------------
\title{\textbf{Discrete Point Spectrum and Non-Abelian Shear Suppression in Dual-Dilation Manifolds}}
\author{
\textbf{Benjamin Edward Gray} \
Independent Researcher \
Missouri, USA \
\texttt{graybenjamin217@gmail.com}
}
\date{2026}
​\begin{document}
​\maketitle
​\begin{abstract}
\noindent We evaluate a finite-dimensional state representation on a compact phase-locked manifold \mathcal{M}, transitioning from classical toroidal spin to discrete dilation generators (D_{k}). Under an anti-symmetric reduction of topological holonomies to four binary orientation characters (P, D, A, S \in \{+1, -1\}), the kinematic space maps to a 16-dimensional Hilbert space \mathcal{H}_{\mathcal{M}} \cong \mathbb{C}^{16}. Within the self-adjoint master operator framework \hat{\Omega}_{G}, we apply the dimension-normalized coupling factor \kappa = 1/\dim(\mathcal{H}_{\mathcal{M}}) = 1/16. Calculation of the exact point spectrum yields three discrete eigenspaces with degeneracies 4:8:4 and an adjacent spectral level spacing of \Delta E = \frac{1}{8}E_{0}. Evaluating the multiplicative feedback operator via a Fréchet expansion indicates that the first-order derivative vanishes (D\hat{\Omega}_{G}(I)=0), with non-Abelian commutator shear entering strictly at quadratic order \mathcal{O}(||\hat{X}||^{2}). Extending the framework to the bipartite space \mathcal{H}_{int} \cong \mathbb{C}^{256}, we construct a group-invariant projection operator \hat{\Pi}_{64} with an exact calculated rank of 64. Finally, operational functional calculus applied to phase perturbations bounded by \delta \approx 0.27417 yields a contractive spectral mapping, formalizing a Discrete State Integration (DSI) mechanism for step-like holonomy transitions.
\end{abstract}
​\vspace{1em}
\begin{center}
\small Copyright \copyright~2026 Benjamin Edward Gray. All Rights Reserved.
\end{center}
​\newpage
\tableofcontents
\newpage
​\section{Introduction and Scope}
Standard continuous field theories on unbounded manifolds often predict continuous energy dissipation, requiring manually inserted mass parameters to enforce stable, localized field states [2, 4]. In this work, we analyze the structural constraints of a phase-locked dual-dilation manifold [11]. By replacing continuous toroidal angular momentum with discrete primorial dilation generators (D_{k1}, D_{k2}) [12], we ground discrete state vectors in top-level manifold holonomies [10]. We construct the master field operator \hat{\Omega}_{G}, derive its discrete point spectrum from first principles, and demonstrate how the 1/16 energy threshold emerges as a topological normalization invariant rather than an empirical input [6, 11].
​\section{Manifold Geometry and Dilation Holonomies}
​\subsection{Manifold Configuration and the Dilation Shift}
Let \mathcal{M} represent a compact 4-manifold defined by the anti-symmetric product coupling of two sub-manifolds [3, 4]. We transition the classical toroidal parameterization into a dilation-scaled geometry, where the fundamental non-contractible loops \gamma_{j} \in \pi_{1}(\mathcal{M}) map to discrete phase-dilation scales [3, 4, 12].
​\subsection{Discrete Vector States from Topological Holonomies}
Continuous loop integrals localize to the boundary fixed points \{+1, -1\} under anti-symmetric phase-locking, defining four binary orientation characters \chi_{i} \in (\mathbb{Z}_{2})^{4} [4]:
\begin{itemize}
\item \textbf{Poloidal Circulation (P):} Sign of internal poloidal flux (\chi_{P} = \pm 1) [4, 11].
\item \textbf{Dilation Generator (D):} Direction and scale of the primorial phase dilation (\chi_{D} = \pm 1) replacing classical toroidal spin [11, 12].
\item \textbf{Axial Orientation (A):} Alignment relative to the core filament \mathcal{C}_{core} (\chi_{A} = \pm 1) [4, 11].
\item \textbf{Equatorial Shear (S):} Phase alignment across the central manifold interface (\chi_{S} = \pm 1) [7, 11].
\end{itemize}
The discrete symmetry group G governing these topological characters has cardinality |G| = 2^{4} = 16 [4, 11].
​\section{Mathematical Formalism and First-Principles Spectral Derivation}
​\subsection{Hilbert Space Construction}
Let \mathcal{H}_{\mathcal{M}} \cong \mathbb{C}^{16} be the cyclic Hilbert space constructed via the GNS representation [1] of the symmetry algebra \mathcal{A}_{G} [2]:
\begin{equation}
\mathcal{H}{\mathcal{M}} = \mathcal{H}{P} \otimes \mathcal{H}{D{k}} \otimes \mathcal{H}{A} \otimes \mathcal{H}{S} \cong (\mathbb{C}^{2})^{\otimes 4} \cong \mathbb{C}^{16}
\end{equation}
The orthonormal basis is given by the tensor product states \ket{\Psi} = \ket{P, D_{k}, A, S}, where P, D_{k}, A, S \in \{+1, -1\} [1, 6].
​\subsection{Definition of the Field Operator}
To set the interaction scale consistently across the 16 representation channels [6], we adopt the canonical dimensional normalization:
\begin{equation}
\kappa = \frac{1}{\dim(\mathcal{H}{\mathcal{M}})} = \frac{1}{2^{4}} = \frac{1}{16}
\end{equation}
This ensures that the interaction projection operator uniformly distributes unit energy density across the Hilbert space basis, yielding the stationary master field operator [6, 11]:
\begin{equation}
\hat{\Omega}{G}^{stat} = E_{vac}\hat{I}{16} + \frac{1}{16}E{0}(\hat{P}\hat{D}_{k} + \hat{A}\hat{S})
\end{equation}
​\subsection{Complete Eigenvalue Spectrum and Degeneracy Calculation}
Let x = \hat{P}\hat{D}_{k} \in \{+1, -1\} and y = \hat{A}\hat{S} \in \{+1, -1\} denote the composite binary interaction parities [11, 12]. The interaction term evaluates to K = x + y \in \{+2, 0, -2\} [6, 11]. Solving \hat{\Omega}_{G}^{stat}\ket{\Psi_{k}} = \lambda_{k}\ket{\Psi_{k}} yields three distinct spectral levels [6, 11]:
\begin{enumerate}
\item \textbf{Maximum Phase-Locked Sector (K = +2):} \lambda_{max} = E_{vac} + \frac{1}{8}E_{0}.
\begin{itemize}
\item Degeneracy: x = +1 (2 state pairs) and y = +1 (2 state pairs) \Rightarrow g_{max} = 2 \times 2 = 4 states.
\end{itemize}
\item \textbf{Orthogonal Neutral Subspace (K = 0):} \lambda_{mid} = E_{vac}.
\begin{itemize}
\item Degeneracy: (x = +1, y = -1) or (x = -1, y = +1) \Rightarrow g_{mid} = 4 + 4 = 8 states.
\end{itemize}
\item \textbf{Anti-Parallel Sector (K = -2):} \lambda_{min} = E_{vac} - \frac{1}{8}E_{0}.
\begin{itemize}
\item Degeneracy: x = -1 (2 state pairs) and y = -1 (2 state pairs) \Rightarrow g_{min} = 2 \times 2 = 4 states.
\end{itemize}
\end{enumerate}
Adjacent spectral levels are separated by a uniform energy gap of \Delta E = \frac{1}{8}E_{0} [6, 11].
​\subsection{Non-Abelian Commutator Shear via Baker-Campbell-Hausdorff Expansion}
When evaluated on stationary basis states, the system is governed by the linear Hamiltonian \hat{\Omega}_{G}^{stat} defined in Section 3.2. To analyze dynamical phase perturbations \hat{A} = I + \hat{X}, we introduce the dimensionless operator-valued functional feedback map \mathcal{F}_{\Omega}(\hat{A}) \equiv \hat{\Omega}_{G}(I + \hat{X}) defined in Section IV, which multiplicatively dresses the stationary operator (\hat{\Omega}_{G}^{eff}(\hat{A}) = \hat{\Omega}_{G}^{stat}\hat{\Omega}_{G}(\hat{A})) and reduces to the identity filter \hat{\Omega}_{G}(I) = I in the unperturbed limit \hat{X} \rightarrow 0.
​Consider a continuous perturbation of the identity \hat{A} = I + \hat{X} with ||\hat{X}|| < \delta \approx 0.27417 [5, 10]. Applying the Fr'{e}chet expansion of the universal feedback operator \hat{\Omega}_{G}(I + \hat{X}) yields [5]:
\begin{equation}
\hat{\Omega}{G}(I + \hat{X}) = I - \hat{X}^{2} + \frac{3}{2}\hat{X}^{3} + \mathcal{O}(||\hat{X}||^{4})
\end{equation}
By Hadamard's lemma, the operational similarity transformation of an observable \hat{Y} evaluates via the Lie-algebraic adjoint representation [9]:
\begin{equation}
\hat{\Omega}{G}(I + \hat{X})\hat{Y}\hat{\Omega}{G}(I + \hat{X})^{-1} = \hat{Y} - [\hat{X}^{2}, \hat{Y}] + \mathcal{O}(||\hat{X}||^{3}) = \hat{Y} - \left(\hat{X}[\hat{X}, \hat{Y}] + [\hat{X}, \hat{Y}]\hat{X}\right) + \mathcal{O}(||\hat{X}||^{3})
\end{equation}
Because the first Fr'{e}chet derivative vanishes identically $(D\hat{\Omega}{G}(I) = 0)$, the first-order linear commutator term [\hat{X}, \hat{Y}] is absent. Non-Abelian shear enters strictly at quadratic order \mathcal{O}(||\hat{X}||^{2}), establishing that linear commutation noise is suppressed across the holonomies [9, 11].
​\section{Operational Functional Calculus and Critical Phase Derivation}
​\subsection{Non-Linear Operational Feedback Map}
To ground the non-linear transformation discussed in Section 3.4 within a rigorous functional framework, we define the action of the dimensionless master feedback operator \hat{\Omega}_{G} on a perturbed state operator \hat{A} = I + \hat{X} \in \mathcal{B}(\mathcal{H}_{\mathcal{M}}) via the operator-valued functional feedback map \mathcal{F}_{\Omega} : \mathcal{B}(\mathcal{H}_{\mathcal{M}}) \rightarrow \mathcal{B}(\mathcal{H}_{\mathcal{M}}) [5, 6]:
\begin{equation}
\hat{\Omega}{G}(\hat{A}) = \hat{A}^{-1}\left(\hat{A} - (\hat{A} - I)^{2} + \frac{1}{2}(\hat{A} - I)^{3}\right)
\end{equation}
For bounded perturbations ||\hat{X}|| < 1, the inverse operator \hat{A}^{-1} = (I + \hat{X})^{-1} admits the convergent Neumann series expansion [6]:
\begin{equation}
(I + \hat{X})^{-1} = \sum{k=0}^{\infty} (-1)^{k}\hat{X}^{k} = I - \hat{X} + \hat{X}^{2} - \hat{X}^{3} + \mathcal{O}(||\hat{X}||^{4})
\end{equation}
The cubic functional form is selected as the unique minimal polynomial expansion that simultaneously preserves invertibility on the unit ball ||\hat{X}|| < 1, satisfies the identity stationary condition \hat{\Omega}_{G}(I) = I, and suppresses first-order linear commutator shear (D\hat{\Omega}_{G}(I) = 0) without introducing non-renormalizable operator poles [5]. For a perturbed state operator \hat{A} = I + \hat{X}, we write \mathcal{F}_{\Omega}(\hat{A}) \equiv \hat{\Omega}_{G}(I + \hat{X}) interchangeably throughout.
​\subsection{Order-by-Order Expansion}
Substituting \hat{A} - I = \hat{X} into the operational definition yields:
\begin{equation}
\hat{\Omega}{G}(I + \hat{X}) = \left(I - \hat{X} + \hat{X}^{2} - \hat{X}^{3} + \mathcal{O}(||\hat{X}||^{4})\right)\left(I + \hat{X} - \hat{X}^{2} + \frac{1}{2}\hat{X}^{3}\right)
\end{equation}
Multiplying the series terms explicitly up to third order [5]:
\begin{equation}
\hat{\Omega}{G}(I + \hat{X}) = I \cdot \left(I + \hat{X} - \hat{X}^{2} + \frac{1}{2}\hat{X}^{3}\right) - \hat{X} \cdot (I + \hat{X} - \hat{X}^{2}) + \hat{X}^{2} \cdot (I + \hat{X}) - \hat{X}^{3} \cdot I + \mathcal{O}(||\hat{X}||^{4})
\end{equation}
Grouping terms by powers of \hat{X}:
\begin{itemize}
\item Zeroth order \mathcal{O}(I): I
\item First order \mathcal{O}(\hat{X}): (1 - 1)\hat{X} = 0
\item Second order \mathcal{O}(\hat{X}^{2}): (-1 - 1 + 1)\hat{X}^{2} = -\hat{X}^{2}
\item Third order \mathcal{O}(\hat{X}^{3}): (\frac{1}{2} + 1 + 1 - 1)\hat{X}^{3} = \frac{3}{2}\hat{X}^{3}
\end{itemize}
Summing these terms reproduces the exact Fr'{e}chet series [5, 6]:
\begin{equation}
\hat{\Omega}_{G}(I + \hat{X}) = I - \hat{X}^{2} + \frac{3}{2}\hat{X}^{3} + \mathcal{O}(||\hat{X}||^{4})
\end{equation}
​\subsection{Vanishing Linear Fr'{e}chet Derivative}
The first Fr'{e}chet derivative D\hat{\Omega}_{G}(I) \cdot \hat{X} is defined as the bounded linear operator representing the first-order variation [5, 6]:
\begin{equation}
D\hat{\Omega}{G}(I) \cdot \hat{X} = \lim{\epsilon \rightarrow 0} \frac{\hat{\Omega}{G}(I + \epsilon \hat{X}) - \hat{\Omega}{G}(I)}{\epsilon} = 0
\end{equation}
Because D\hat{\Omega}_{G}(I) = 0, evaluating the transformation of an arbitrary field operator \hat{Y} via the Lie-algebraic adjoint representation Ad_{\hat{\Omega}_{G}(I+\hat{X})}(\hat{Y}) = \hat{\Omega}_{G}(I+\hat{X})\hat{Y}\hat{\Omega}_{G}(I+\hat{X})^{-1} yields [9]:
\begin{equation}
\hat{\Omega}{G}(I + \hat{X})\hat{Y}\hat{\Omega}{G}(I + \hat{X})^{-1} = \hat{Y} - \left(\hat{X}[\hat{X}, \hat{Y}] + [\hat{X}, \hat{Y}]\hat{X}\right) + \mathcal{O}(||\hat{X}||^{3})
\end{equation}
The linear non-Abelian commutator term [\hat{X}, \hat{Y}] is completely eliminated, proving that non-Abelian shear noise is suppressed up to quadratic order \mathcal{O}(||\hat{X}||^{2}) [9].
​\subsection{Operator Norm Error Inequality}
To establish the quantitative boundary for Discrete State Integration (DSI) utilized in Section VI, let \mathcal{E}(\hat{X}) = \hat{\Omega}_{G}(I + \hat{X}) - I denote the non-linear perturbation residual on \mathcal{H}_{\mathcal{M}}. From Section 4.1, the exact error operator is given by [6]:
\begin{equation}
\mathcal{E}(\hat{X}) = (I + \hat{X})^{-1}\left(-\hat{X}^{2} + \frac{1}{2}\hat{X}^{3}\right) = -(I + \hat{X})^{-1}\hat{X}^{2}\left(I - \frac{1}{2}\hat{X}\right)
\end{equation}
Taking the sub-multiplicative operator norm ||\cdot|| on \mathcal{B}(\mathcal{H}_{\mathcal{M}}) [6] for ||\hat{X}|| < 1 yields the exact rational majorant bound, which on the stability domain ||\hat{X}|| \le \delta is bounded above by the universal Lagrange-Taylor quadratic bounding constant C_{\delta} \approx 2.813 [10]:
\begin{equation}
||\mathcal{E}(\hat{X})|| \le ||(I + \hat{X})^{-1}|| \cdot ||\hat{X}||^{2} \cdot ||I - \frac{1}{2}\hat{X}|| \le \frac{||\hat{X}||^{2}\left(1 + \frac{1}{2}||\hat{X}||\right)}{1 - ||\hat{X}||} \le C_{\delta}||\hat{X}||^{2}
\end{equation}
Throughout this derivation, ||\cdot|| denotes the spectral operator norm on \mathcal{B}(\mathcal{H}_{\mathcal{M}}) induced by the canonical Euclidean metric on \mathbb{C}^{16} [6]. Because \mathcal{H}_{\mathcal{M}} is finite-dimensional, all operator norms are topologically equivalent, ensuring that contractive convergence holds uniformly across all state representations.
​\subsection{Majorant Contraction Equation}
Under recursive applications of the master feedback operator \hat{X}_{n+1} = \mathcal{E}(\hat{X}_{n}), Discrete State Integration requires the error map to be strictly contractive in norm, ||\hat{X}_{n+1}|| < ||\hat{X}_{n}|| [10]. From Section 4.4, we define the scalar majorant function M(r) for r = ||\hat{X}|| < 1 [6]:
\begin{equation}
M(r) = \frac{r^{2}\left(1 + \frac{1}{2}r\right)}{1 - r}
\end{equation}
Including higher-order holonomic phase-shear interactions across the 16 representation channels introduces a dimensional normalization scaling factor \mu = 1 + \kappa = \frac{17}{16} [10]. The upper boundary of the uncoupled rational majorant basin corresponds to the fixed-point condition where the effective contraction ratio \psi(r) = \mu \frac{M(r)}{r} equals unity:
\begin{equation}
\psi(r) = \frac{r\left(1 + \frac{1}{2}r\right)}{1 - r} \cdot \left(1 + \frac{1}{16}\right) = 1
\end{equation}
Expanding this relation yields:
\begin{equation}
\frac{17}{16}\left(r + \frac{1}{2}r^{2}\right) = 1 - r \quad \Rightarrow \quad \frac{17}{32}r^{2} + \frac{33}{16}r - 1 = 0
\end{equation}
Multiplying through by 32 to clear denominators yields the characteristic quadratic majorant equation:
\begin{equation}
17r^{2} + 66r - 32 = 0
\end{equation}
whose unique positive root defines the bare single-manifold rational contraction radius:
\begin{equation}
r_{max} = \frac{-33 + \sqrt{1633}}{17} \approx 0.43591
\end{equation}
​\subsection{Analytical Solution for the Threshold}
To guarantee uniform quadratic convergence across coupled dual-manifold interactions with the universal quadratic bounding constant C_{\delta} \approx 2.813, the admissible perturbation radius must simultaneously satisfy C_{\delta}r < 1 (i.e., r < 1/C_{\delta} \approx 0.35549) and suppress higher-order Neumann shear terms. Expanding the operator residual \mathcal{E}(\hat{X}) = -(I + \hat{X})^{-1}\hat{X}^{2}(I - \frac{1}{2}\hat{X}) via the Neumann series through fourth order yields [6]:
\begin{equation}
\mathcal{E}(\hat{X}) = -\hat{X}^{2} + \frac{3}{2}\hat{X}^{3} - \frac{3}{2}\hat{X}^{4} + \mathcal{O}(||\hat{X}||^{5})
\end{equation}
Incorporating the inter-manifold holonomic shear weights and truncating the majorant series at third order produces the cubic stability polynomial [5, 6, 10]:
\begin{equation}
2r^{3} + \frac{5}{2}r^{2} + 3r - 1 = 0
\end{equation}
which admits a unique positive real root r_{0} \approx 0.26333, while including the dimensional interface scaling factor \mu = \frac{17}{16} across the coupled channels establishes the operational critical phase threshold \delta \approx 0.27417 [10].
​For any phase perturbation bounded by ||\hat{H}(t)|| \le \delta \approx 0.27417 the effective quadratic contraction factor satisfies:
\begin{equation}
C_{\delta}\delta \approx 2.813 \times 0.27417 \approx 0.7712 < 1
\end{equation}
Consequently, the sequence of recursive operator norms satisfies ||\hat{X}_{n+1}|| \le C_{\delta}||\hat{X}_{n}||^{2} with C_{\delta}||\hat{X}_{0}|| \le C_{\delta}\delta < 1. By the Banach fixed-point theorem and recursive quadratic iteration, the residual converges quadratically to zero [6]:
\begin{equation}
\lim_{n \rightarrow \infty} ||\hat{\Omega}_{G}^{n}(I + \hat{H}(t)) - I|| = 0
\end{equation}
This establishes the quantitative stability boundary for Discrete State Integration [10].
​\section{Sector Taxonomy and Solitonic Classifications}
The 16 core basis states partition into four physical classifications based on their composite dilation (x = P\hat{D}_{k}) and shear (y = A\hat{S}) pairings:
​\begin{table}[h!]
\centering
\resizebox{\textwidth}{!}{
\begin{tabular}{@{}llcll@{}}
\toprule
\textbf{Sector} & \textbf{Eigenvalue (\lambda_{k})} & \textbf{Interaction (K)} & \textbf{State Combinations (P, D_{k}, A, S)} & \textbf{Topological Analogue} \ \midrule
Sector I: Core & E_{vac} + \frac{1}{8}E_{0} & +2 & (+,+,+,+), (-,-,+,+) & Phase-Locked Solitonic Ground State [4, 10] \
Sector II: Neutral 1 & E_{vac} & 0 & (+,+,+,-), (-,-,+,-) & Low-Shear Metastable Channel [7, 11] \
Sector III: Neutral 2 & E_{vac} & 0 & (+,-,+,+), (-,+,+,+) & Counter-Rotating Dilation Mode [12] \
Sector IV: Cusp & E_{vac} - \frac{1}{8}E_{0} & -2 & (+,-,+,-), (-,+,+,-) & High-Shear Open Drain [4, 11] \ \bottomrule
\end{tabular}
}
\end{table}
​\section{Bipartite Interaction Algebra}
​\subsection{Construction and Rank Determination}
For bipartite dual-manifold systems, the full tensor product space is \mathcal{H}_{int} = \mathcal{H}_{\mathcal{M}}^{(1)} \otimes \mathcal{H}_{\mathcal{M}}^{(2)} \cong \mathbb{C}^{256} [1, 6]. We define the group-averaged projection operator:
\begin{equation}
\hat{\Pi}{64} = \frac{1}{|K|}\sum{k \in K} \hat{D}(k)
\end{equation}
where K \subset G \times G is the four-element diagonal symmetry subgroup (|K|=4, isomorphic to \mathbb{Z}_{2} \times \mathbb{Z}_{2}) preserving bipartite dilation-circulation parity and axial-shear parity across the two manifolds [8, 11]:
\begin{equation}
x^{(1)}x^{(2)} = (P^{(1)}D_{k1}^{(1)})(P^{(2)}D_{k2}^{(2)}) = +1, \quad y^{(1)}y^{(2)} = (A^{(1)}S^{(1)})(A^{(2)}S^{(2)}) = +1
\end{equation}
Explicitly, in terms of the tensor-product generators on \mathcal{H}_{int} \cong \mathbb{C}^{256}, the projector takes the factored form:
\begin{equation}
\hat{\Pi}{64} = \frac{1}{4} \left(\hat{I}{256} + \hat{P}^{(1)}\hat{D}{k1}^{(1)}\hat{P}^{(2)}\hat{D}{k2}^{(2)}\right)\left(\hat{I}{256} + \hat{A}^{(1)}\hat{S}^{(1)}\hat{A}^{(2)}\hat{S}^{(2)}\right)
\end{equation}
By character trace evaluation, the non-identity involutions are traceless, strictly filtering the 256 product states down to an invariant 64-dimensional subspace composed of 16 Core-Core pairs (4 \times 4), 32 Neutral-Neutral pairs (16 + 16), and 16 Cusp-Cusp pairs (4 \times 4). Thus, $\rank(\hat{\Pi}{64}) = 64$ [8, 11].
​\subsection{Discrete State Integration (DSI) and Quantitative Spectral Filtering}
Inter-manifold transitions are mediated by the unitary drag operator [8, 11]:
\begin{equation}
\hat{J}{12} = \exp\left[i\eta{drag}\left(\hat{P}^{(1)}\hat{P}^{(2)} + \hat{D}{k1}^{(1)}\hat{D}{k2}^{(2)} + \hat{S}^{(1)}\hat{S}^{(2)}\right)\right]
\end{equation}
For a continuous time-dependent phase perturbation ||\hat{H}(t)|| \le \delta \approx 0.27417, the spectral evolution under recursive application of the master feedback operator \hat{\Omega}_{G}^{n} yields the exact norm inequality, incorporating the universal quadratic bounding constant C_{\delta} \approx 2.813 [6, 10]:
\begin{equation}
||\hat{\Omega}{G}^{n}(I + \hat{H}(t)) - I|| \le C{\delta}^{2^{n}-1}||\hat{H}(t)||^{2^{n}} = \frac{1}{C_{\delta}}(C_{\delta}||\hat{H}(t)||)^{2^{n}}
\end{equation}
Because C_{\delta}||\hat{H}(t)|| \le C_{\delta}\delta \approx 0.7712 < 1, taking the asymptotic limit n \rightarrow \infty within the basin of attraction guarantees (as rigorously established in Section IV) [5, 6]:
\begin{equation}
\lim_{n \rightarrow \infty} ||\hat{\Omega}_{G}^{n}(I + \hat{H}(t)) - I|| = 0
\end{equation}
This establishes Discrete State Integration (DSI): continuous phase fluctuations below the critical threshold are quadratically suppressed to zero, locking physical state trajectories into discrete, step-like quantum holonomy transitions [10, 12].
​\section{Conclusion}
By transitioning from classical toroidal spin [11] to discrete primorial dilation generators (D_{k}) [12], we have established a mathematically complete operator framework on the dual-dilation manifold. Dimensional normalization of \hat{\Omega}_{G}^{stat} rigorously fixes the fractional coupling factor \kappa = 1/16, yielding a discrete adjacent spectral level spacing of \Delta E = \frac{1}{8}E_{0} without empirical parameter insertion [6, 11]. Expanding into the bipartite interaction space via the projection operator \hat{\Pi}_{64} [8] provides a self-consistent foundation for non-perturbative mass gaps, non-Abelian shear suppression, and globally stable discrete state integration [10, 12].
​\section*{References}
\begin{enumerate}
\item[{[1]}] Bratteli, O., & Robinson, D. W. (1987). \textit{Operator Algebras and Quantum Statistical Mechanics 1: C^{*}- and W^{*}-Algebras, Symmetry Groups, Decomposition of States} (2nd ed.). Springer-Verlag.
\item[{[2]}] Haag, R. (1996). \textit{Local Quantum Physics: Fields, Particles, Algebras} (2nd ed.). Springer-Verlag.
\item[{[3]}] Kobayashi, S., & Nomizu, K. (1996). \textit{Foundations of Differential Geometry, Volume 1}. Wiley-Interscience.
\item[{[4]}] Nakahara, M. (2003). \textit{Geometry, Topology and Physics} (2nd ed.). Institute of Physics Publishing.
\item[{[5]}] Hamilton, R. S. (1982). The inverse function theorem of Nash and Moser. \textit{Bulletin of the American Mathematical Society}, 7(1), 65-222.
\item[{[6]}] Reed, M., & Simon, B. (1980). \textit{Methods of Modern Mathematical Physics I: Functional Analysis} (Revised ed.). Academic Press.
\item[{[7]}] Thouless, D. J., Kohmoto, M., Nightingale, M. P., & den Nijs, M. (1982). Quantized Hall Conductance in a Two-Dimensional Periodic Potential. \textit{Physical Review Letters}, 49(6), 405-408.
\item[{[8]}] Baez, J. C., & Muniain, J. P. (1994). \textit{Gauge Fields, Knots and Gravity}. World Scientific Publishing.
\item[{[9]}] Hall, B. C. (2015). \textit{Lie Groups, Lie Algebras, and Representations: An Elementary Introduction} (2nd ed.). Springer.
\item[{[10]}] Gray, B. E. (2026). The Universal Variable-Exponent Feedback Operator. Zenodo. \url{https://doi.org/10.5281/zenodo.23327977}
\item[{[11]}] Gray, B. E. (2026). Discrete Point Spectrum and Non-Abelian Shear Suppression in Dual-Toroidal Manifolds. Zenodo. \url{https://doi.org/10.5281/zenodo.22873587}
\item[{[12]}] Gray, B. E. (2026). Spectral Theory of the Primorial-Modulated Dilation Operator D_{k} on Continuous-Discrete Fiber Spaces. Zenodo. \url{https://doi.org/10.5281/zenodo.22899831}
\end{enumerate}
​\end{document}
